\honest{My Best and Worst Mistake}{Eliezer Yudkowsky}{2008}

When I was young, I refused to give ``intelligence'' a mathematical definition, and I was ``loath to put any definition on `intelligence' at all.'' The reason was a standard bait-and-switch in AI. You define ``intelligence'' as, say, ``logical reasoning,'' build a cheap theorem-prover, and say, ``Lo, I have implemented intelligence!'' Such definitions picked out correlates of intelligence instead of its core, and researchers then chased what they had written down. I was not building a career in AI. I wanted a mind that could actually build nanotechnology, so I had no reason to redefine intelligence to puff up a paper. I knew intuitively what I was after, something powerful enough to take stars apart for raw material, and I did not want definitions to distract me from it. \nb{Yudkowsky's archived early texts do define intelligence, though not mathematically. ``Staring into the Singularity'' has a section headed ``The Definition of Smartness.'' ``The Plan to Singularity'' (1999--2000) says: ``Intelligence is whatever lets you model, predict, and manipulate reality.''}

Many of my mistakes came from reacting too far against other people's mistakes, and refusing to define intelligence was one. Another: I had seen AI projects brought down by ``physics envy,'' sticking to elegant math and so to toy systems. I concluded that any math neat enough to fit an equation would not work for real intelligence, ``Except for Bayes's Theorem.'' That exception either softens the offense or shows I should have suspected the whole generalization. I had read too few AI books, too popular, and the wrong ones. I believed the cliché that AI overpromised and had written it off as a ``sick, dead field,'' so I never looked hard enough to find the math that was not fake. \nb{The Plan says that ``AI remains crippled'' and that the needed talent is more likely to be found in cognitive science.}

This disbelief was one of my best mistakes. Expecting no simple answer, I read cognitive psychology, functional and computational neuroanatomy, evolutionary psychology and biology, and several branches of AI. When I had a bright idea I did not rush to implement it: intelligence was a puzzle with many pieces, and a mind missing even one might do nothing interesting. I was wrong that academic AI was a wasteland, and more wrong that there could be no math of intelligence. But I do not regret the neuroanatomy, even though I now think an AI should look nothing like a human brain. It taught me that a mind's parts are things like ``visual cortex'' and ``cerebellum,'' not a ``commonsense reasoning module,'' which is a standard wrong road in AI.

Set aside the wrong reasons and look at what the beliefs made me do. Writing AI off sent me to the cognitive sciences. Expecting no simple answer kept me gathering information instead of stopping at one brilliant idea. Refusing to define intelligence meant I studied the problem for years before proposing any systematization. My moral: ``What you actually end up doing, screens off the clever reason why you're doing it.'' Clever reasoning that sends you to study many sciences leaves you far better off, once it proves stupid, than clever reasoning that says you need not read the books.

Many of my successes were right actions for wrong reasons. You should put that down to the anthropic principle: if I had hit a dead end you probably would not be hearing from me. To me it remains an embarrassment. My Traditional Rationalist upbringing pushed the accidents in a good direction, toward reasons to study, and helped me recover from mistakes; but none of it was the right action for the right reason. I write partly to leave a trail to where I ended up by accident, so that others need less luck.

It was also one of my worst mistakes, because ``informal'' can mean ``held to low standards.'' My clever reason for not defining ``intelligence'' and some other terms, that others had gone astray by defining it, let sloppy reasoning in. Defining intelligence at once would have been wrong too; you cannot define fire before you know about atoms. Until then you are better off saying ``that orangey-bright thing,'' and you have to talk about it to investigate it. But reasoning at that level is something you do on the way to knowing better. You do not put your weight down on it or draw firm conclusions from it, however inescapable it seems. The young Eliezer put his weight down on it and stepped onto a loaded trap. To be continued. \nb{``A Prodigy of Refutation,'' two posts on, returns to the young author's belief that a superintelligence would know what is right.}
